\chapter{Literaturwerte}

In der folgenden Tabelle werden Literaturwerte gesammelt.

\begin{table}[caption={Literaturwerte},label={literaturwerte},]
    Name & Formelzeichen & Literaturwert & Quelle \\
    Adiabatenkoeffizient Kohlendioxid & $\kappa_C$ & 1.3 & Versuchsanleitung\\
    Adiabatenkoeffizient Luft & $\kappa_L$ & 1.4 & Versuchsanleitung\\
    Avogadrokonstante & $N_A$ & 6.02214076 $\cdot 10^{23}$ 1/mol & Wikipedia\\
    Boltzmannkonstante & $k$  & $k = 1.38 \cdot 10^{-23}~J/K$ & Wikipedia \\
    Elementarladung & e & 1.602176634 $\cdot 10^{-19}$ C & Wikipedia \\
    Erdbeschleunigung & g & $(9.80984 \pm 2 \cdot 10^{-5})~m/s^2$ & Versuchsanleitung \\
    Faraday-Konstante & F & 9.64853321233100184 $\cdot 10^4$ C/mol& Wikipedia\\
    ideale Gaskonstante & R & 8.31446261815324 J/mol K & Wikipedia\\
    Lichtgeschwindigkeit & c & 299792458 m/s & Wikipedia \\
    molare Masse von Kupfer & $M_{Mol}^{Cu}$ & 63.5460 g/mol & \href{chemie.de}{chemie.de}\\
    Molekülmasse Kohlendioxid & $M_C$ & 44 g/mol & Versuchsanleitung\\
    Molekülmasse Luft & $M_L$ & 29 g/mol & Versuchsanleitung\\
    Molvolumen ideales Gas (Idealbedingungen) & $V_0^{mol}$ & 22.414 l/mol & Versuchsanleitung\\
    Normalbedingung Druck & $p_0$ & 1013.25 mbar & Versuchsanleitung\\
    Normalbedingung Temperatur & $T_0$ & 273.15 K& Versuchsanleitung\\
    Rydbergkonstante & $R_{\infty}$ & $1.0973731568160(21) \cdot 10^7$ & Wikipedia \\
    Wertigkeit von Kupferionen & $z_{\ce{Cu^{++}}}$ & 2 & Versuchsanleitung\\
    Wertigkeit von Sulfationen & $z_{\ce{SO_4^{--}}}$ & 2 & Versuchsanleitung\\
    Wertigkeit von Sauerstoffmolekülen & $z_{\ce{O_2}}$ & 4 & Versuchsanleitung\\
\end{table}

Der Wert der Erdbeschleunigung bezieht sich auf die Stadt Heidelberg.

Dabei ist zu beachten, dass der Begriff der Wertigkeit bei geladenen Ionen von der chemischen Reaktionsgleichung abhängt. 

\vfill  